\documentclass[10pt,a4paper]{amsart}
\usepackage{amsmath,amssymb,amsthm,graphicx}
\usepackage{verbatim}
\newenvironment{reference}{\comment}{\endcomment}
%\newenvironment{reference}{}{}
\newenvironment{slogan}{\comment}{\endcomment}
\newenvironment{history}{\comment}{\endcomment}

% For commutative diagrams we use Xy-pic
\usepackage[all]{xy}

% We use 2cell for 2-commutative diagrams.
\xyoption{2cell}
% Theorem environments.
%
\theoremstyle{plain}
\newtheorem{theorem}[subsection]{Theorem}
\newtheorem{proposition}[subsection]{Proposition}
\newtheorem{lemma}[subsection]{Lemma}

\theoremstyle{definition}
\newtheorem{definition}[subsection]{Definition}
\newtheorem{example}[subsection]{Example}
\newtheorem{exercise}[subsection]{Exercise}
\newtheorem{situation}[subsection]{Situation}

\theoremstyle{remark}
\newtheorem{remark}[subsection]{Remark}
\newtheorem{remarks}[subsection]{Remarks}

\numberwithin{equation}{subsection}

% Macros
%
\def\lim{\mathop{\mathrm{lim}}\nolimits}
\def\colim{\mathop{\mathrm{colim}}\nolimits}
\def\Spec{\mathop{\mathrm{Spec}}}
\def\Hom{\mathop{\mathrm{Hom}}\nolimits}
\def\Ext{\mathop{\mathrm{Ext}}\nolimits}
\def\SheafHom{\mathop{\mathcal{H}\!\mathit{om}}\nolimits}
\def\SheafExt{\mathop{\mathcal{E}\!\mathit{xt}}\nolimits}
\def\Sch{\mathit{Sch}}
\def\Mor{\mathop{\mathrm{Mor}}\nolimits}
\def\Ob{\mathop{\mathrm{Ob}}\nolimits}
\def\Sh{\mathop{\mathit{Sh}}\nolimits}
\def\NL{\mathop{N\!L}\nolimits}
\def\CH{\mathop{\mathrm{CH}}\nolimits}
\def\proetale{{pro\text{-}\acute{e}tale}}
\def\etale{{\acute{e}tale}}
\def\QCoh{\mathit{QCoh}}
\def\Ker{\mathop{\mathrm{Ker}}}
\def\Im{\mathop{\mathrm{Im}}}
\def\Coker{\mathop{\mathrm{Coker}}}
\def\Coim{\mathop{\mathrm{Coim}}}

% Boxtimes
%
\DeclareMathSymbol{\boxtimes}{\mathbin}{AMSa}{"02}

%
% Macros for moduli stacks/spaces
%
\def\QCohstack{\mathcal{QC}\!\mathit{oh}}
\def\Cohstack{\mathcal{C}\!\mathit{oh}}
\def\Spacesstack{\mathcal{S}\!\mathit{paces}}
\def\Quotfunctor{\mathrm{Quot}}
\def\Hilbfunctor{\mathrm{Hilb}}
\def\Curvesstack{\mathcal{C}\!\mathit{urves}}
\def\Polarizedstack{\mathcal{P}\!\mathit{olarized}}
\def\Complexesstack{\mathcal{C}\!\mathit{omplexes}}
% \Pic is the operator that assigns to X its picard group, usage \Pic(X)
% \Picardstack_{X/B} denotes the Picard stack of X over B
% \Picardfunctor_{X/B} denotes the Picard functor of X over B
\def\Pic{\mathop{\mathrm{Pic}}\nolimits}
\def\Picardstack{\mathcal{P}\!\mathit{ic}}
\def\Picardfunctor{\mathrm{Pic}}
\def\Deformationcategory{\mathcal{D}\!\mathit{ef}}
\usepackage{cleveref}
\crefname{originallemma}{original lemma}{original lemmas}
\Crefname{originallemma}{Original lemma}{Original lemmas}
\makeatletter
\newlabel{algebra-lemma-Artin-Rees}{{00IN}{source}{}{sourcepointer-1}{}}
\newlabel{algebra-lemma-Artin-Rees@cref}{{[originallemma][0][]00IN}{[1][0][]source}{}{}{}}
\makeatother
\crefname{originallemma}{original lemma}{original lemmas}
\Crefname{originallemma}{Original lemma}{Original lemmas}
\makeatletter
\newlabel{algebra-lemma-quotient-complete}{{031A}{source}{}{sourcepointer-2}{}}
\newlabel{algebra-lemma-quotient-complete@cref}{{[originallemma][0][]031A}{[1][0][]source}{}{}{}}
\makeatother
\crefname{originalequation}{original equation}{original equations}
\Crefname{originalequation}{Original equation}{Original equations}
\makeatletter
\newlabel{equation-C-prime}{{0AL4}{source}{}{sourcepointer-3}{}}
\newlabel{equation-C-prime@cref}{{[originalequation][0][]0AL4}{[1][0][]source}{}{}{}}
\makeatother
\crefname{originalreference}{original reference}{original references}
\Crefname{originalreference}{Original reference}{Original references}
\makeatletter
\newlabel{lemma-get-morphism-nonzerodivisor}{{0AK6}{source}{}{sourcepointer-4}{}}
\newlabel{lemma-get-morphism-nonzerodivisor@cref}{{[originalreference][0][]0AK6}{[1][0][]source}{}{}{}}
\makeatother
\crefname{originalreference}{original reference}{original references}
\Crefname{originalreference}{Original reference}{Original references}
\makeatletter
\newlabel{lemma-topologically-finite-type-Noetherian}{{0AJQ}{source}{}{sourcepointer-5}{}}
\newlabel{lemma-topologically-finite-type-Noetherian@cref}{{[originalreference][0][]0AJQ}{[1][0][]source}{}{}{}}
\makeatother
\title{Lifting maps through powers of a principal ideal modulo the a-power torsion}
\author{The Stacks Project Authors}
\date{stacks-project:restricted.tex:local-snapshot-20261012; GNU FDL 1.2 or later}
\begin{document}
\maketitle
\noindent\textit{Editorial note.} All delivered mathematical text is the author's own, copied from the pinned Stacks Project chapter file restricted.tex; no statement or proof step was rewritten, reordered or removed.\par\medskip
\noindent\textit{Editorial note.} The counted proof is the author's complete argument: the proof quotients by the torsion submodule, verifies completeness of the quotient, uses the exact rows to view C as a fibre product, lifts to the quotient by the nonzerodivisor lemma, and reassembles the map via the fibre product description\par\medskip
\noindent\textit{Editorial note.} The label lemma-get-morphism-principal names the counted statement, and the proof environment immediately following it in the source is delivered verbatim as its proof.\par\medskip
\noindent\textit{Editorial note.} The counted statement is carried by the author's own lemma environment in the chapter source.\par\medskip
\noindent\textit{Allowed standard prerequisites (editorial).}\begin{itemize}
\item the lifting lemma for principal ideals with nonzerodivisor (lemma-get-morphism-nonzerodivisor of the same chapter)
\item completeness of quotients: C/C[a\textasciicircum{}infty] is I-adically complete when the intersection of C[a\textasciicircum{}infty] + a\textasciicircum{}nC equals C[a\textasciicircum{}infty] (Algebra, Lemma algebra-lemma-quotient-complete)
\item the fibre product description of C from the exact sequence with the torsion submodule
\item existence of the integer d with C[a\textasciicircum{}infty] intersect a\textasciicircum{}dC = 0: for Noetherian C the torsion stabilizes at C[a\textasciicircum{}e] and Artin-Rees gives the intersection bound (Algebra, Lemma algebra-lemma-Artin-Rees)
\end{itemize}\medskip
\ifdefined\hypertarget\hypertarget{sourcepointer-1}{}\fi
\noindent\textit{Source reference.} Original lemma 00IN: sources/stacks-chow/tags.txt, source lines 685--685. Prerequisite pointer invoked at proof line(s) 1309; it names a standard fact of a sibling chapter that the counted argument uses, and every step of the counted proof built on it is retained verbatim in the delivered body.\par\medskip
\ifdefined\hypertarget\hypertarget{sourcepointer-2}{}\fi
\noindent\textit{Source reference.} Original lemma 031A: sources/stacks-chow/tags.txt, source lines 3977--3977. Prerequisite pointer invoked at proof line(s) 1319; it names a standard fact of a sibling chapter that the counted argument uses, and every step of the counted proof built on it is retained verbatim in the delivered body.\par\medskip
\ifdefined\hypertarget\hypertarget{sourcepointer-3}{}\fi
\noindent\textit{Source reference.} Original equation 0AL4: sources/stacks-restricted/pinned-original-source.tex, source lines 326--326. Prerequisite pointer invoked at proof line(s) 1296; it names a standard fact of the same chapter that the counted argument uses, and every step of the counted proof built on it is retained verbatim in the delivered body.\par\medskip
\ifdefined\hypertarget\hypertarget{sourcepointer-4}{}\fi
\noindent\textit{Source reference.} Original reference 0AK6: sources/stacks-restricted/pinned-original-source.tex, source lines 1242--1242. Prerequisite pointer invoked at proof line(s) 1335; it names a standard fact of the same chapter that the counted argument uses, and every step of the counted proof built on it is retained verbatim in the delivered body.\par\medskip
\ifdefined\hypertarget\hypertarget{sourcepointer-5}{}\fi
\noindent\textit{Source reference.} Original reference 0AJQ: sources/stacks-restricted/pinned-original-source.tex, source lines 395--395. Prerequisite pointer invoked at proof line(s) 1314; it names a standard fact of the same chapter that the counted argument uses, and every step of the counted proof built on it is retained verbatim in the delivered body.\par\medskip
\begin{lemma}
\label{lemma-get-morphism-principal}
Let $I = (a)$ be a principal ideal of a Noetherian ring $A$.
Let $B$ be an object of (\ref{equation-C-prime}).
Assume given an integer $c \geq 0$ such that $\Ext^1_B(\NL_{B/A}^\wedge, N)$
is annihilated by $a^c$ for all $B$-modules $N$.
Let $C$ be an $I$-adically complete $A$-algebra.
Assume given an integer $d \geq 0$ such that $C[a^\infty] \cap a^dC = 0$.
Let $n > \max(2c, c + d)$. For any $A$-algebra map
$\psi_n : B \to C/a^nC$ there exists an $A$-algebra map
$\varphi : B \to C$ such
that $\psi_n \bmod a^{n - c} = \varphi \bmod a^{n - c}$.
\end{lemma}

\begin{proof}
Let $C \to C'$ be the quotient of $C$ by $C[a^\infty]$. The $A$-algebra
$C'$ is $I$-adically complete by
Algebra, Lemma \ref{algebra-lemma-quotient-complete}
and the fact that $\bigcap (C[a^\infty] + a^nC) = C[a^\infty]$
because for $n \geq d$ the sum $C[a^\infty] + a^nC$ is direct.
For $m \geq d$ the diagram
$$
\xymatrix{
0 \ar[r] &
C[a^\infty] \ar[r] \ar[d] &
C \ar[r] \ar[d] & C' \ar[r] \ar[d] & 0 \\
0 \ar[r] &
C[a^\infty] \ar[r] &
C/a^m C \ar[r] & C'/a^m C' \ar[r] & 0
}
$$
has exact rows. Thus $C$ is the fibre product of $C'$ and
$C/a^mC$ over $C'/a^mC'$ for all $m \geq d$. By
Lemma \ref{lemma-get-morphism-nonzerodivisor}
we can choose a homomorphism $\varphi' : B \to C'$
such that $\varphi'$ and $\psi_n$ agree as maps into $C'/a^{n - c}C'$.
We obtain a homomorphism $(\varphi', \psi_n \bmod a^{n - c}C) : B \to
C' \times_{C'/a^{n - c}C'} C/a^{n - c}C$.
Since $n - c \geq d$ this is the same thing as a homomorphism
$\varphi : B \to C$. This finishes the proof.
\end{proof}

\end{document}
